% GENERATED FILE. Edit canonical lessons, then run npm run generate:books.
% canonical-source-hash: fnv1a64:8cef4c73d7a7ed9f
% canonical-lessons: UR-C301-sahafi, UR-C301-mauqa, UR-C301-editar, UR-C301-tafsil, UR-C301-likhvana

\chapter{Reporters}
\label{ch:ur-reporters}
\hlchaptermodality{301}

\begin{chapteropening}
\textbf{By the end of this chapter:} \emph{I can say who reported a story and from where.}

Complete the last lesson of chapter 301: I can say who reported a story and from where.
\end{chapteropening}

\section[\texorpdfstring{\ur{صحافی} (sahāfī)}{صحافی (sahāfī)}]{\ur{صحافی} (sahāfī) — a journalist}

\label{lesson:UR-C301-sahafi}



\begin{quote}
Before the new one: say the Urdu for an official notice, then the Urdu for a sensation.
\end{quote}

\subsection*{You'll want to know: \ur{صحافی}}
\textbf{\ur{صحافی}} — \emph{sahāfī} — \textquotedblleft{}a journalist\textquotedblright{}.

\begin{quote}
\ur{صحافی} \ur{نے} \ur{داروغے} \ur{سے} \ur{سوال} \ur{پوچھے۔}
\end{quote}

\begin{quote}
\emph{sahāfī ne dāroghe se savāl pūchhe.}
\end{quote}

\begin{quote}
The journalist asked the inspector questions.
\end{quote}

Masculine or feminine, by the person meant.

\begin{cousinweb}
\textbf{\ur{صحافی}} (\emph{sahāfī}) is usually traced, through Persian, to Arabic \emph{ṣaḥāfa}, \textquotedblleft{}the press, journalism\textquotedblright{}, with the ending \emph{-ī}: one who writes for the papers.
\end{cousinweb}

\begin{grammarlens}[title={What you've built}]
One more word for telling what happened.
\end{grammarlens}

\subsection*{Your turn}
\begin{itemize}
\raggedright
  \item \emph{Say it:} \emph{sahāfī}
  \item \emph{Say it:} \emph{sahāfī}, once more
  \item \emph{Say it:} \emph{sahāfī ne dāroghe se savāl pūchhe.} — The journalist asked the inspector questions.
  \item \emph{Recall:} say the Urdu for an official notice, then the Urdu for a sensation, then say \emph{sahāfī} again
\end{itemize}

\begin{tcolorbox}[breakable,colback=teal!4,colframe=teal!35!black,title={Before you move on}]
What does \ur{صحافی} mean? (\textquotedblleft{}A journalist\textquotedblright{}.) Say it once more.
\end{tcolorbox}
\section[\texorpdfstring{\ur{موقع} (mauqa)}{موقع (mauqa)}]{\ur{موقع} (mauqa) — the scene of an event; an occasion}

\label{lesson:UR-C301-mauqa}



\begin{quote}
Before the new one: say the Urdu for a sensation, then the Urdu for a journalist.
\end{quote}

\subsection*{You'll want to know: \ur{موقع}}
\textbf{\ur{موقع}} — \emph{mauqa} — \textquotedblleft{}the scene of an event; an occasion\textquotedblright{}.

\begin{quote}
\ur{صحافی} \ur{موقع} \ur{پر} \ur{پہنچا۔}
\end{quote}

\begin{quote}
\emph{sahāfī mauqe par pahuṅchā.}
\end{quote}

\begin{quote}
The journalist reached the scene.
\end{quote}

Masculine.

\begin{cousinweb}
\textbf{\ur{موقع}} (\emph{mauqa}) is usually traced, through Persian, to Arabic \emph{mawqiʿ}, \textquotedblleft{}a place where something falls, a site\textquotedblright{}. \textbf{\ur{موقع} \ur{پر}} is on the spot; the word also means a chance or an occasion.
\end{cousinweb}

\begin{grammarlens}[title={What you've built}]
One more word for telling what happened.
\end{grammarlens}

\subsection*{Your turn}
\begin{itemize}
\raggedright
  \item \emph{Say it:} \emph{mauqa}
  \item \emph{Say it:} \emph{mauqa}, once more
  \item \emph{Say it:} \emph{sahāfī mauqe par pahuṅchā.} — The journalist reached the scene.
  \item \emph{Recall:} say the Urdu for a sensation, then the Urdu for a journalist, then say \emph{mauqa} again
\end{itemize}

\begin{tcolorbox}[breakable,colback=teal!4,colframe=teal!35!black,title={Before you move on}]
What does \ur{موقع} mean? (\textquotedblleft{}The scene of an event; an occasion\textquotedblright{}.) Say it once more.
\end{tcolorbox}
\section[\texorpdfstring{\ur{ایڈیٹر} (eḍīṭar)}{ایڈیٹر (eḍīṭar)}]{\ur{ایڈیٹر} (eḍīṭar) — an editor}

\label{lesson:UR-C301-editar}



\begin{quote}
Before the new one: say the Urdu for a journalist, then the Urdu for the scene of an event.
\end{quote}

\subsection*{You'll want to know: \ur{ایڈیٹر}}
\textbf{\ur{ایڈیٹر}} — \emph{eḍīṭar} — \textquotedblleft{}an editor\textquotedblright{}.

\begin{quote}
\ur{ایڈیٹر} \ur{نے} \ur{خبر} \ur{چھاپی۔}
\end{quote}

\begin{quote}
\emph{eḍīṭar ne khabar chhāpī.}
\end{quote}

\begin{quote}
The editor printed the story.
\end{quote}

Masculine or feminine, by the person meant.

\begin{cousinweb}
is borrowed from English \emph{editor}. The \textbf{\ur{ایڈیٹر}} decides what the paper prints.
\end{cousinweb}

\begin{grammarlens}[title={What you've built}]
One more word for telling what happened.
\end{grammarlens}

\subsection*{Your turn}
\begin{itemize}
\raggedright
  \item \emph{Say it:} \emph{eḍīṭar}
  \item \emph{Say it:} \emph{eḍīṭar}, once more
  \item \emph{Say it:} \emph{eḍīṭar ne khabar chhāpī.} — The editor printed the story.
  \item \emph{Recall:} say the Urdu for a journalist, then the Urdu for the scene of an event, then say \emph{eḍīṭar} again
\end{itemize}

\begin{tcolorbox}[breakable,colback=teal!4,colframe=teal!35!black,title={Before you move on}]
What does \ur{ایڈیٹر} mean? (\textquotedblleft{}An editor\textquotedblright{}.) Say it once more.
\end{tcolorbox}
\section[\texorpdfstring{\ur{تفصیل} (tafsīl)}{تفصیل (tafsīl)}]{\ur{تفصیل} (tafsīl) — a detail; a full account}

\label{lesson:UR-C301-tafsil}



\begin{quote}
Before the new one: say the Urdu for the scene of an event, then the Urdu for an editor.
\end{quote}

\subsection*{You'll want to know: \ur{تفصیل}}
\textbf{\ur{تفصیل}} — \emph{tafsīl} — \textquotedblleft{}a detail; a full account\textquotedblright{}.

\begin{quote}
\ur{اخبار} \ur{میں} \ur{حادثے} \ur{کی} \ur{تفصیل} \ur{تھی۔}
\end{quote}

\begin{quote}
\emph{akhbār meṅ hādse kī tafsīl thī.}
\end{quote}

\begin{quote}
The paper had the details of the accident.
\end{quote}

Feminine.

\begin{cousinweb}
\textbf{\ur{تفصیل}} (\emph{tafsīl}) is usually traced, through Persian, to Arabic \emph{tafṣīl}, \textquotedblleft{}a setting out part by part\textquotedblright{}. \textbf{\ur{تفصیل} \ur{سے}} is in detail.
\end{cousinweb}

\begin{grammarlens}[title={What you've built}]
One more word for telling what happened.
\end{grammarlens}

\subsection*{Your turn}
\begin{itemize}
\raggedright
  \item \emph{Say it:} \emph{tafsīl}
  \item \emph{Say it:} \emph{tafsīl}, once more
  \item \emph{Say it:} \emph{akhbār meṅ hādse kī tafsīl thī.} — The paper had the details of the accident.
  \item \emph{Recall:} say the Urdu for the scene of an event, then the Urdu for an editor, then say \emph{tafsīl} again
\end{itemize}

\begin{tcolorbox}[breakable,colback=teal!4,colframe=teal!35!black,title={Before you move on}]
What does \ur{تفصیل} mean? (\textquotedblleft{}A detail; a full account\textquotedblright{}.) Say it once more.
\end{tcolorbox}
\section[\texorpdfstring{\ur{لکھوانا} (likhvānā)}{لکھوانا (likhvānā)}]{\ur{لکھوانا} (likhvānā) — to have written down; to dictate}

\label{lesson:UR-C301-likhvana}



\begin{quote}
Before the new one: say the Urdu for an editor, then the Urdu for a detail.
\end{quote}

\subsection*{You'll want to know: \ur{لکھوانا}}
\textbf{\ur{لکھوانا}} — \emph{likhvānā} — \textquotedblleft{}to have written down; to dictate\textquotedblright{}.

\begin{quote}
\ur{گواہ} \ur{نے} \ur{اپنا} \ur{بیان} \ur{لکھوایا۔}
\end{quote}

\begin{quote}
\emph{gavāh ne apnā bayān likhvāyā.}
\end{quote}

\begin{quote}
The witness had his statement written down.
\end{quote}

Past, with \ur{نے} (he or she …): \ur{اس} \ur{نے} \ur{رپورٹ} \ur{لکھوائی}, \emph{us ne riporṭ likhvāī}.

\begin{cousinweb}
\textbf{\ur{لکھوانا}} (\emph{likhvānā}) is the second causing form of \textbf{\ur{لکھنا}}, to write, which CDIAL lists under Sanskrit \emph{likháti}, \textquotedblleft{}scratches, writes\textquotedblright{} (CDIAL 11048). The ending \emph{-vā-} hands the writing to someone else.
\end{cousinweb}

\begin{grammarlens}[title={What you've built}]
One more word for telling what happened.
\end{grammarlens}

\subsection*{Your turn}
\begin{itemize}
\raggedright
  \item \emph{Say it:} \emph{likhvānā}
  \item \emph{Say it:} \emph{likhvānā}, once more
  \item \emph{Say it:} \emph{gavāh ne apnā bayān likhvāyā.} — The witness had his statement written down.
  \item \emph{Recall:} say the Urdu for an editor, then the Urdu for a detail, then say \emph{likhvānā} again
\end{itemize}

\begin{tcolorbox}[breakable,colback=teal!4,colframe=teal!35!black,title={Before you move on}]
What does \ur{لکھوانا} mean? (\textquotedblleft{}To have written down; to dictate\textquotedblright{}.) Say it once more.
\end{tcolorbox}
